%%%%%%%%%%%%%%%%%%%%%%%%%%%%%%%%%%%%%%%%%%%%%%%%%%%%%%%%%%%%%%%
%
% Welcome to Overleaf --- just edit your LaTeX on the left,
% and we'll compile it for you on the right. If you open the
% 'Share' menu, you can invite other users to edit at the same
% time. See www.overleaf.com/learn for more info. Enjoy!
%
%%%%%%%%%%%%%%%%%%%%%%%%%%%%%%%%%%%%%%%%%%%%%%%%%%%%%%%%%%%%%%%
\documentclass{article}
\usepackage{amsmath, amssymb}
\begin{document}
Soient $f$ et $g$ $\in \mathbb{F}[x]$ deux polynômes non nuls de degrés respectifs
n et m, avec $n$,$m$ $\in \mathbb{N}$, on a :
\[
\begin{aligned}
    f &= f_nx^n+\cdots+f_1x+f_0 \\
    g &= g_nx^m+\cdots+g_1x+g_0
\end{aligned}
\]
\\
Le résultant de f et g est le déterminant de l'application linéaire 
\[
\varphi :
\mathbb{F}[x]_{<m} \times \mathbb{F}[x]_{<n}
\longrightarrow
\mathbb{F}[x]_{<m+n},
\qquad
(s,t) \longrightarrow sf + tg,
\]
\\
La matrice de l'application linéaire $\varphi$ est appelée la matrice de Sylvester associée à $f$ et $g$ et est notée $S(f,g)$.
\\
Le résultant de $f$ et $g$ est alors donné par :
\[
\operatorname{Res}(f,g)=\det S(f,g)
\]
La matrice de Sylvester $S(f,g)$ s'écrit sous la forme suivante :
\[
S(f,g)=
\begin{pmatrix}
f_n &        &        &        & g_m &        &        \\
f_{n-1} & f_n &        &        & g_{m-1} & g_m &        \\
\vdots & \vdots & \ddots &        & \vdots & \vdots & \ddots \\
f_0 & f_1 & \cdots & f_n & g_0 & g_1 & \cdots \\
    & f_0 & \ddots & \vdots &     & g_0 & \ddots \\
    &     & \ddots & f_1 &     &     & \ddots \\
    &     &        & f_0 &     &     & g_0
\end{pmatrix}
\]
Les sous-résultants de $f$ et $g$ sont définis à partir de la matrice de
Sylvester $S(f,g)$.
Pour tout entier $k \ge 0$, le $k$-ième sous-résultant de $f$ et $g$,
noté $\mathrm{SRes}_k(f,g)$, est le déterminant d’une sous-matrice
appropriée de $S(f,g)$.
La suite des sous-résultants permet de déterminer le degré du
$\gcd(f,g)$ : si
\[
\mathrm{SRes}_0(f,g)=\cdots=\mathrm{SRes}_{d-1}(f,g)=0
\quad \text{et} \quad
\mathrm{SRes}_d(f,g)\neq 0,
\]
alors $\deg(\gcd(f,g))=d$, et $\mathrm{SRes}_d(f,g)$ est proportionnel
au $\gcd(f,g)$.


\end{document}